\subsection{Stochastic Escape and Noise-Sensitive Switching}

Real systems are noisy. Fatigue, interruption, novelty, and mood fluctuations
all behave like state disturbances. This matters because barrier crossing can
occur either through deliberate control or through stochastic escape.

\begin{discussion}
The stochastic viewpoint helps prevent overmoralization. A person may drift out
of a good work state not because they positively chose failure, but because the
state was shallow and noise-sensitive. Conversely, a good intervention may
work not because it increases virtue, but because it deepens the desired basin
enough that ordinary noise no longer knocks the trajectory out.
\end{discussion}
